% ======================================================================
% Chapter 19: Perturbation Theory Anomalies
% (book pages 651-686; PDF pages 674-711)
% ======================================================================
\chapter{Perturbation Theory Anomalies}

In many examples, we have seen that loop corrections can have an important
effect on the predictions of quantum field theory.  We have studied examples
in which the relative importance of operators is shifted by radiative
corrections, and in which the form of the interactions they mediate is
altered.  However, in specific circumstances, radiative corrections can have an
even more significant effect: They can destroy symmetries of the classical
equations of motion.

The most important effect of this type involves the chiral symmetries of
theories with massless fermions.  In Section~3.4, we saw that the massless
Dirac Lagrangian has an enhanced symmetry associated with the separate number
conservation of left- and right-handed fermions.  This symmetry is generated
by the axial vector current $j^{\mu 5} = \bar{\psi}\gamma^\mu\gamma^5\psi$.
Classically,
\begin{equation}
  \partial_\mu j^{\mu 5} = 0
  \label{eq:19.1}
\end{equation}
for zero-mass fermions.  This equation of motion is true not only in free
fermion theory but also, as a classical field equation, in massless QED and
QCD.

However, in this chapter, we will see that the true picture is not so simple.
We will show that, in gauge theories, the conservation of the axial vector
current is actually incompatible with gauge invariance, and that radiative
corrections in gauge theories supply a nonzero operator that appears on the
right-hand side of Eq.~\eqref{eq:19.1}.  This new conservation equation for
the axial current has a number of remarkable consequences, which we will
discuss in Sections~19.3 and~19.4.

\section{The Axial Current in Two Dimensions}

Eventually, we will want to analyze the current conservation equation for the
axial current in massless QCD.  However, this discussion will involve some
technical complication, so we will first study the physics that violates axial
current conservation in a context in which the calculations are relatively
simple.  A particularly simple model problem is that of two-dimensional
massless QED.

The Lagrangian of two-dimensional QED is
\begin{equation}
  \mathcal{L} = \bar{\psi}(i\slashed{D})\psi - \frac{1}{4}(F_{\mu\nu})^2 ,
  \label{eq:19.2}
\end{equation}
with $\mu, \nu = 0, 1$ and $D_\mu = \partial_\mu + ieA_\mu$.  The Dirac matrices
must be chosen to satisfy the Dirac algebra
\begin{equation}
  \{\gamma^\mu, \gamma^\nu\} = 2g^{\mu\nu} .
  \label{eq:19.3}
\end{equation}
In two dimensions, this set of relations can be represented by $2\times2$
matrices; we choose
\begin{equation}
  \gamma^0 = \sigma^2,\qquad \gamma^1 = i\sigma^1 .
  \label{eq:19.4}
\end{equation}
The Dirac spinors will be two-component fields.

The product of the Dirac matrices, which anticommutes with each of the
$\gamma^\mu$, is
\begin{equation}
  \gamma^5 = \gamma^0\gamma^1 = \sigma^3 .
  \label{eq:19.5}
\end{equation}
Then, just as in four dimensions, there are two possible currents,
\begin{equation}
  j^\mu = \bar{\psi}\gamma^\mu\psi;\qquad j^{\mu 5} = \bar{\psi}\gamma^\mu
    \gamma^5\psi ,
  \label{eq:19.6}
\end{equation}
and both are conserved if there is no mass term in the Lagrangian.

To make the conservation laws quite explicit, we label the components of the
fermion field $\psi$ in this spinor basis as
\begin{equation}
  \psi = \begin{pmatrix}\psi_+\\ \psi_-\end{pmatrix} .
  \label{eq:19.7}
\end{equation}
The subscript indicates the $\gamma^5$ eigenvalue.  Then, using the explicit
representations \eqref{eq:19.4} and \eqref{eq:19.7}, we can rewrite the
fermionic part of \eqref{eq:19.2} as
\begin{equation}
  \mathcal{L} = \psi^{\dagger}_+ i(D_0 + D_1)\psi_+
    + \psi^{\dagger}_- i(D_0 - D_1)\psi_- .
  \label{eq:19.8}
\end{equation}
In the free theory, the field equation of $\psi_+$ would be
\begin{equation}
  i(\partial_0 + \partial_1)\psi_+ = 0 ;
  \label{eq:19.9}
\end{equation}
the solutions to this equation are waves that move to the right in the
one-dimensional space at the speed of light.  We will thus refer to the
particles associated with $\psi_+$ as \emph{right-moving fermions}.  The
quanta associated with $\psi_-$ are, similarly, \emph{left-moving}.  This
distinction is analogous to the distinction between left- and right-handed
particles which gives the physical interpretation of $\gamma^5$ in four
dimensions.  Since the Lagrangian \eqref{eq:19.8} contains no terms that mix
left- and right-moving fields, it seems obvious that the number currents for
these fields are separately conserved.  Thus,
\begin{equation}
  \partial_0 N_+ = \partial_0 N_- = 0 .
  \label{eq:19.10}
\end{equation}

It is a curious property of two-dimensional spacetime that the vector and
axial vector fermionic currents are not independent of each other.  Let
$\epsilon^{\mu\nu}$ be the totally antisymmetric symbol in two dimensions,
with $\epsilon^{01} = +1$.  Then the two-dimensional Dirac matrices obey the
identity
\begin{equation}
  \gamma^\mu\gamma^5 = -\epsilon^{\mu\nu}\gamma_\nu .
  \label{eq:19.11}
\end{equation}
The currents $j^{\mu 5}$ and $j^\mu$ have the same relation.  Thus we can study
the properties of the axial vector current by using results that we have
already derived for the vector current.

\subsection{Vacuum Polarization Diagrams}

In Section~7.5, we computed the lowest-order vacuum polarization of QED in
dimensional regularization.  In the limit of zero mass, we found, in
Eq.~(7.90),
\begin{equation}
  i\Pi^{\mu\nu}(q) = -i(q^2g^{\mu\nu} - q^\mu q^\nu)\ \frac{2e^2}{(4\pi)^{d/2}}\
    \mathrm{tr}[\mathbb{1}]\ \int_0^1 dx\ x(1-x)\ \frac{\Gamma(2-d/2)}
      {(-x(1-x)q^2)^{2-d/2}} ,
  \label{eq:19.12}
\end{equation}
where $\mathrm{tr}[\mathbb{1}] = 4$ gives the convention for tracing over
Dirac matrices given in Eq.~(7.88).  If we set
$\mathrm{tr}[\mathbb{1}] = 2$ to be consistent with \eqref{eq:19.4} and then
set $d = 2$ in \eqref{eq:19.12}, we find the finite and well-defined result
\begin{equation}
  \Pi^{\mu\nu}(q) = -\frac{e^2}{\pi}\ \frac{q^\mu q^\nu - q^2g^{\mu\nu}}
    {q^2} .
  \label{eq:19.13}
\end{equation}
Notice that this expression has the structure of a photon mass term; the
photon receives the mass
\begin{equation}
  m_\gamma^2 = \frac{e^2}{\pi} .
  \label{eq:19.14}
\end{equation}
Schwinger showed that this result is exact, and that the photon of
two-dimensional QED is a free massive boson.%
\footnote{J.~Schwinger, \emph{Phys.\ Rev.} \textbf{128}, 2425 (1962).}  In the
discussion below Eq.~\eqref{eq:7.72}, we pointed out that it is not possible
for a vacuum polarization amplitude consistent with the Ward identity to
generate a mass for the photon unless it also contains a pole at $q^2 = 0$.
In two dimensions, such a pole can arise from the infrared behavior of the
fermion-antifermion intermediate state, and we see this behavior explicitly in
\eqref{eq:19.13}.

\begin{figure}[t]
  \centering
  \includegraphics[width=0.5\textwidth]{fig19.1.png}
  \caption{Computation of $\langle j^\mu\rangle$ in a background
    electromagnetic field.}
  \label{fig:19.1}
\end{figure}

Once we have an explicit expression for the vacuum polarization, we can find
the expectation value of the current induced by a background electromagnetic
field.  This quantity is generated by the diagram of Fig.~\ref{fig:19.1},
which gives
\begin{equation}
  \int d^2x\ e^{iq\cdot x}\ \langle j^\mu(x)\rangle
  = -\bigl( \Pi^{\mu\nu}(q) \bigr)\ A_\nu(q)
  = -\frac{e^2}{\pi}\ \frac{q^\mu q^\nu - q^2g^{\mu\nu}}{q^2}\ A_\nu(q) .
  \label{eq:19.15}
\end{equation}
where $A_\nu(q)$ is the Fourier transform of the background field.  This
quantity manifestly satisfies the current conservation relation
$q_\mu\langle j^\mu(q)\rangle = 0$.

The identity \eqref{eq:19.11} between the vector and axial vector currents
allows us to derive from \eqref{eq:19.15} the corresponding expectation value
of $j^{\mu 5}$.  We find
\begin{equation}
  \langle j^{\mu 5}(q)\rangle = -\epsilon^{\mu\nu}\ \langle j_\nu(q)\rangle
  = \frac{e^2}{\pi}\ \frac{\epsilon^{\mu\nu}q_\nu q^\lambda}{q^2}\ A_\lambda(q)
  .
  \label{eq:19.16}
\end{equation}
If the axial vector current were conserved, this object would satisfy the Ward
identity.  Instead, we find
\begin{equation}
  q_\mu\langle j^{\mu 5}(q)\rangle = \frac{e^2}{2\pi}\ \epsilon^{\alpha\beta}
    F_{\alpha\beta}(q) .
  \label{eq:19.17}
\end{equation}
This is the Fourier transform of the field equation
\begin{equation}
  \partial_\mu j^{\mu 5} = \frac{e^2}{2\pi}\ \epsilon^{\alpha\beta}F_{\alpha\beta}
  .
  \label{eq:19.18}
\end{equation}
Apparently, the axial vector current is not conserved in the presence of
electromagnetic fields, as the result of an anomalous behavior of its vacuum
polarization diagram.

How could this happen?  The Feynman diagrams formally satisfy the Ward
identity both for the vector and for the axial vector current.  The problem
must come in the regularization of the vacuum polarization diagram.  By
dimensional analysis, we know that this diagram has the form
\begin{equation}
  \Pi^{\mu\nu}(q) = A\ q^\mu q^\nu + B\ g^{\mu\nu}q^2 .
  \label{eq:19.19}
\end{equation}
The coefficient $B$ is a finite integral, and is, in any event, unambiguously
determined by the low-energy structure of the theory since it is the residue
of the pole in $q^2$.  However, the integral $A$ is logarithmically divergent,
so its value depends on the regularization.  Dimensional regularization
automatically subtracts this integral to set $A = -B$; then the vector
current Ward identity is satisfied.  But then we are led directly to
\eqref{eq:19.17}.  We could, alternatively, regularize the integral $A$ so that
$A = 0$.  Working through the steps of the previous paragraph with this
modification, we now find $q_\mu\langle j^{\mu 5}(q)\rangle = 0$, but
\begin{equation}
  q_\mu\langle j^\mu(q)\rangle = \frac{e^2}{\pi}\ q_\mu A^\mu(q) .
  \label{eq:19.20}
\end{equation}
Though the result \eqref{eq:19.17} is unpleasant, the result \eqref{eq:19.20}
would be a complete disaster, since it depends on the unphysical gauge
degrees of freedom of the vector potential.  We conclude that it is not
possible to regularize two-dimensional QED so that, simultaneously, the theory
is gauge invariant and the axial vector current is conserved.  The price of
requiring gauge invariance is the anomalous nonconservation of the axial
current shown in \eqref{eq:19.18}.

\subsection{The Axial Vector Current Operator Equation}

To understand what happened to the axial current from another viewpoint, we
now study the operator equation for the divergence of $j^{\mu 5}$.  Varying
the Lagrangian \eqref{eq:19.2}, we find the following equations of motion for
the fermion fields:
\begin{equation}
  i\slashed{D}\psi = 0,\qquad iD_\mu\bar{\psi}\gamma^\mu = 0 .
  \label{eq:19.21}
\end{equation}
By using these equations of motion in the most straightforward way, it is easy
to conclude that $\partial_\mu j^{\mu 5} = 0$.  However, a closer look at
these manipulations reveals some subtleties, which alter the final conclusion.

The axial vector current is a composite operator built out of fermion
fields.  In the previous chapter we saw that products of local operators are
often singular, so we will define the current by placing the two fermion
fields at distinct points separated by a distance $\epsilon$ and then carefully
taking the limit as the two fields approach each other.  Explicitly, we define
\begin{equation}
  j^{\mu 5}(x) = \mathrm{symm}\lim_{\epsilon\to 0}\ \Bigl[
    \bar{\psi}\Bigl( x + \frac{\epsilon}{2} \Bigr)\gamma^\mu\gamma^5\
    \exp\Bigl( -ie\int_{x-\epsilon/2}^{x+\epsilon/2} dz\cdot A(z) \Bigr)\
    \psi\Bigl( x - \frac{\epsilon}{2} \Bigr) \Bigr] .
  \label{eq:19.22}
\end{equation}
Notice that, because we have placed $\bar{\psi}$ and $\psi$ at different
points, we must introduce a Wilson line \eqref{eq:15.53} in order that the
operator be locally gauge invariant.  To give $j^{\mu 5}$ the correct
transformation properties under Lorentz transformations, the limit
$\epsilon\to 0$ should be taken symmetrically,
\begin{equation}
  \mathrm{symm}\lim_{\epsilon\to 0}\Bigl\{ \frac{\epsilon^\mu}{\epsilon^2}
    \Bigr\} = 0,\qquad
  \mathrm{symm}\lim_{\epsilon\to 0}\Bigl\{ \frac{\epsilon^\mu\epsilon^\nu}
    {\epsilon^2} \Bigr\} = \frac{g^{\mu\nu}}{d} ,
  \label{eq:19.23}
\end{equation}
with $d = 2$ in this case.

We now compute the divergence of the axial current defined as in
\eqref{eq:19.22}.  Using the equations of motion, we obtain the operator
equation
\begin{equation}
  \partial_\mu j^{\mu 5} = \mathrm{symm}\lim_{\epsilon\to 0}\ \Bigl[
    \bar{\psi}\Bigl( x + \frac{\epsilon}{2} \Bigr)\gamma^\mu\gamma^5\
    \exp(\cdots)\ \psi\Bigl( x - \frac{\epsilon}{2} \Bigr) \Bigr]\
    \cdots ,
  \label{eq:19.25}
\end{equation}
where the singular terms in the product of operators at short distances must
be evaluated carefully.  The leading term in the short-distance expansion of
the product $\bar{\psi}(x+\epsilon/2)\psi(x-\epsilon/2)$ is obtained by
contracting the two fields with a free-field propagator.  In two dimensions,
this gives a singular contribution that, when traced with
$\gamma^\mu\gamma^5$, yields a nonzero result in the limit $\epsilon\to 0$.
This is the operator form of the anomaly:
\begin{equation}
  \partial_\mu j^{\mu 5} = \frac{e^2}{2\pi}\ \epsilon^{\alpha\beta}F_{\alpha\beta}
  .
  \label{eq:19.29}
\end{equation}

\subsection{An Example with Fermion Number Nonconservation}

The anomalous nonconservation of the axial current has a dramatic
consequence: it implies that the numbers of left- and right-handed fermions
can be changed by background gauge fields.  To see this in a simple example,
we consider two-dimensional QED on a circle of length $L$, with the fermions
satisfying periodic boundary conditions.  We couple the theory to a constant
background electric field, which corresponds to a time-dependent $A^1$
component.

The Hamiltonian of the system is
\begin{equation}
  H = \int dx\ \Bigl[ \psi^{\dagger}_+(i\partial_1 - eA^1)\psi_+
    + \psi^{\dagger}_-(-i\partial_1 + eA^1)\psi_- \Bigr] ,
  \label{eq:19.34}
\end{equation}
where $\psi_{\pm}$ are the right- and left-moving components.  For a constant
$A^1$ field, it is easy to diagonalize this Hamiltonian.  The eigenstates of
the covariant derivatives are wavefunctions
\begin{equation}
  e^{ik_n x},\qquad \text{with } k_n = \frac{2\pi}{L}n,\qquad n = -\infty,
    \ldots, \infty ,
  \label{eq:19.35}
\end{equation}
to satisfy the periodic boundary conditions.  Then the single-particle
eigenstates of $H$ have energies
\begin{equation}
  E^+_n = +(k_n - eA^1),\qquad E^-_n = -(k_n - eA^1) .
  \label{eq:19.36}
\end{equation}
Each type of fermion has an infinite tower of equally spaced levels.  To find
the ground state of $H$, we fill the negative energy levels and interpret
holes created among these filled states as antiparticles.

\begin{figure}[t]
  \centering
  \includegraphics[width=0.55\textwidth]{fig19.2.png}
  \caption{Effect on the vacuum state of the Hamiltonian $H$ of
    one-dimensional QED due to an adiabatic change in the background $A^1$
    field.}
  \label{fig:19.2}
\end{figure}

Now adiabatically change the value of $A^1$.  The fermion energy levels
slowly shift in accord with the relations \eqref{eq:19.36}.  If $A^1$ changes
by the finite amount
\begin{equation}
  \Delta A^1 = \frac{2\pi}{eL} ,
  \label{eq:19.37}
\end{equation}
which brings the Wilson loop Eq.~(19.32) back to its original value, the
spectrum of $H$ returns to its original form.  In this process, each level of
$\psi_+$ moves down to the next position, and each level of $\psi_-$ moves up
to the next position, as shown in Fig.~\ref{fig:19.2}.  The occupation
numbers of levels should be maintained in this adiabatic process.  Thus,
remarkably, one right-moving fermion disappears from the vacuum and one extra
left-moving fermion appears.  At the same time,
\begin{equation}
  \int d^2x\ \frac{e}{2\pi}\ \epsilon^{\mu\nu}F_{\mu\nu}
  = \int dt\,dx\ \frac{e}{\pi}\ \partial_0 A^1
  = \frac{e}{\pi}\ \Delta A^1\,L
  = \frac{e}{\pi}\cdot\frac{2\pi}{eL}\cdot L = 2 ,
  \label{eq:19.38}
\end{equation}
where we have inserted \eqref{eq:19.37} in the last line.  Thus the integrated
form of the anomalous nonconservation equation \eqref{eq:19.18} is indeed
satisfied:
\begin{equation}
  N_R - N_L = \int d^2x\ \frac{e}{2\pi}\ \epsilon^{\mu\nu}F_{\mu\nu} .
  \label{eq:19.39}
\end{equation}
Even in this simple example, we see that it is not possible to escape the
question of ultraviolet regularization in analyzing the chiral conservation
law.  Right-moving fermions are lost and left-moving fermions appear from the
depths of the fermionic spectrum, $E\to -\infty$.  In computing the changes
in the separate fermion numbers, we have assumed that the vacuum cannot
change the charge it contains at large negative energies.  This prescription
is gauge invariant, but it leads to the nonconservation of the axial vector
current.

\section{The Axial Current in Four Dimensions}

All of the derivations we have just given for the two-dimensional axial
current have analogues in four dimensions.  In Eq.~\eqref{eq:3.40}, we showed
that, in the case of massless fermions, the four-dimensional Dirac equation
splits neatly into separate equations for left- and right-handed fermions.  If
we couple the Dirac equation to a gauge field, we replace derivatives by
covariant derivatives.  This does not seem to affect the manifest separation
between the two helicity components.  Thus it seems clear that both the
vector and axial vector currents should remain conserved.  However, after the
analysis we have just completed for the two-dimensional case, we know that we
should not take these conservation laws for granted.  We will now make a more
careful analysis of the axial vector conservation law in four dimensions.

\subsection{The Axial Vector Current Operator Equation}

We begin with the case of massless four-dimensional QED.  Of the three
arguments that we gave in the previous section for the two-dimensional axial
current conservation law, the operator derivation generalizes most easily.
The fermion field equations \eqref{eq:19.21} are identical in the
four-dimensional case.  We can again adopt the gauge-invariant definition of
the axial vector current \eqref{eq:19.22}.  When we take the divergence of
this current, all of the manipulations leading to Eq.~\eqref{eq:19.25} are
still correct.

\begin{figure}[t]
  \centering
  \includegraphics[width=0.5\textwidth]{fig19.3.png}
  \caption{Expansion of $\bar{\psi}(y)\psi(z)$ in the presence of a
    background gauge field.}
  \label{fig:19.3}
\end{figure}

From this point, we must compute the singular terms in the operator product
of the two fermion fields in the limit $\epsilon\to 0$.  As in two dimensions,
the leading term is given by contracting the two operators using a free-field
propagator.  This contribution gives
\begin{equation}
  \bigl\langle \bar{\psi}(y)\psi(z) \bigr\rangle \to -i\slashed{\partial}\
    \frac{1}{2\pi^2(y-z)^2} = -i\gamma^\alpha\ \frac{(y-z)_\alpha}
      {2\pi^2(y-z)^4} .
  \label{eq:19.40}
\end{equation}
This is highly singular as $(y-z)\to 0$, but it gives zero when traced with
$\gamma^\mu\gamma^5$.  To find a nonzero result, we must consider terms of
higher order in the expansion of the product of operators.

In a nonzero background gauge field, the contraction of fermion fields is
given by the series of diagrams shown in Fig.~\ref{fig:19.3}.  We have
computed the leading term in this series in \eqref{eq:19.40}.  The higher
terms give less singular terms as $(y-z)\to 0$.  The second term in the
series is given by
\begin{equation}
  \text{(second term)} = \int \frac{d^4k}{(2\pi)^4}\ \frac{d^4p}{(2\pi)^4}\
    e^{ik\cdot y}\ e^{-ip\cdot z}\ \frac{(-ie\slashed{A}(p))}{k^2(k+p)^2}
    \ \cdots ,
  \label{eq:19.41}
\end{equation}
This contribution leads to
\begin{equation}
  \Bigl\langle \bar{\psi}\Bigl( x+\frac{\epsilon}{2} \Bigr)\gamma^\mu\gamma^5
    \psi\Bigl( x-\frac{\epsilon}{2} \Bigr) \Bigr\rangle
  \sim -4e\ \epsilon^{\mu\alpha\beta\gamma}\ \partial_\beta
    \partial_\gamma A_\alpha(x)\ \Bigl( \frac{1}{\epsilon^2} \Bigr)\ \cdots .
  \label{eq:19.42}
\end{equation}
To evaluate the limit $\epsilon\to 0$, we can expand the integrand for large
$k$.  Then
\begin{equation}
  \Bigl\langle \bar{\psi}\Bigl( x+\frac{\epsilon}{2} \Bigr)\gamma^\mu\gamma^5
    \psi\Bigl( x-\frac{\epsilon}{2} \Bigr) \Bigr\rangle
  \to 4e\ \epsilon^{\mu\nu\rho\sigma}\ \partial_\rho A_\sigma(x)\ \Bigl(
    \frac{1}{4\pi^2}\log\frac{1}{\epsilon^2} \Bigr) .
  \label{eq:19.43}
\end{equation}
Substituting this expression into \eqref{eq:19.25}, we find
\begin{equation}
  \partial_\mu j^{\mu 5} = \mathrm{symm}\lim_{\epsilon\to 0}\ \Bigl[
    \frac{e^2}{8\pi^2}\ \epsilon^{\mu\nu\rho\sigma}F_{\mu\nu}F_{\rho\sigma}
    \Bigr] .
  \label{eq:19.44}
\end{equation}
Now take the symmetric limit $\epsilon\to 0$ in four dimensions.  We find
\begin{equation}
  \partial_\mu j^{\mu 5} = \frac{e^2}{16\pi^2}\ \epsilon^{\mu\nu\rho\sigma}
    F_{\mu\nu}F_{\rho\sigma} .
  \label{eq:19.45}
\end{equation}
This is the \emph{Adler-Bell-Jackiw anomaly} equation.  The right-hand side
is a new operator that appears in the divergence of the axial current as a
result of radiative corrections.

\subsection{Chiral Transformation of the Functional Integral}

It is interesting to derive the axial anomaly in yet another way, using the
functional integral.  This derivation is particularly transparent and
generalizes easily to other theories.  Consider the functional integral for
massless QED,
\begin{equation}
  Z = \int \mathcal{D}\bar{\psi}\,\mathcal{D}\psi\,\mathcal{D}A\ \exp\Bigl[
    i\int d^4x\ \bigl( \bar{\psi}i\slashed{D}\psi
      - \frac{1}{4}(F_{\mu\nu})^2 \bigr) \Bigr] .
  \label{eq:19.61}
\end{equation}
We now perform a chiral transformation on the fermion fields,
\begin{equation}
  \psi(x) \to e^{i\alpha(x)\gamma^5}\psi(x),\qquad
  \bar{\psi}(x) \to \bar{\psi}(x)e^{i\alpha(x)\gamma^5} .
  \label{eq:19.62}
\end{equation}
The Lagrangian transforms as
\begin{equation}
  \bar{\psi}i\slashed{D}\psi \to \bar{\psi}i\slashed{D}\psi
    - \bar{\psi}\slashed{D}\alpha\,\gamma^5\psi
    = \bar{\psi}i\slashed{D}\psi + \alpha\,\partial_\mu j^{\mu 5} .
  \label{eq:19.63}
\end{equation}
If the fermionic measure $\mathcal{D}\bar{\psi}\,\mathcal{D}\psi$ were
invariant under the chiral transformation, then varying the transformed
integral with respect to $\alpha(x)$ would give $\partial_\mu j^{\mu 5} = 0$,
the naive conservation law.  However, we will now show that the measure is
not invariant, and that the Jacobian of the transformation produces exactly
the anomalous term.

To compute the Jacobian, we expand the fermion field in a basis of
eigenstates of $i\slashed{D}$.  Define right and left eigenvectors of
$i\slashed{D}$ by
\begin{equation}
  i\slashed{D}\phi_m = \lambda_m\phi_m,\qquad
  \phi^{\dagger}_m\,i\slashed{D} = \lambda^{*}_m\phi^{\dagger}_m .
  \label{eq:19.64}
\end{equation}
For zero background $A^\mu$ field, these eigenstates are Dirac wavefunctions
of definite momentum; the eigenvalues satisfy
\begin{equation}
  \lambda^2_m = k^2 = (k^0)^2 - (\mathbf{k})^2 .
  \label{eq:19.65}
\end{equation}
For a fixed $A^\mu$ field, this is also the asymptotic form of the eigenvalues
for large $k$.  These eigenfunctions give us a basis that we can use to expand
$\psi$ and $\bar{\psi}$:
\begin{equation}
  \psi(x) = \sum_m a_m\phi_m(x),\qquad
  \bar{\psi}(x) = \sum_m \bar{a}_m\phi^{\dagger}_m(x) ,
  \label{eq:19.66}
\end{equation}
where $a_m$, $\bar{a}_m$ are anticommuting coefficients multiplying the
c-number eigenfunctions \eqref{eq:19.64}.  The functional measure over
$\psi$, $\bar{\psi}$ can then be defined as
\begin{equation}
  \mathcal{D}\bar{\psi}\,\mathcal{D}\psi = \prod_m d\bar{a}_m\,da_m ,
  \label{eq:19.67}
\end{equation}
and the functional measure over $\psi'$, $\bar{\psi}'$ can be defined in the
same way.

If $\psi'(x) = (1+i\alpha(x)\gamma^5)\psi(x)$, the expansion coefficients of
$\psi$ and $\psi'$ are related by an infinitesimal linear transformation
$(1+C)$, computed as follows:
\begin{equation}
  a'_m = \sum_n \int d^4x\ \phi^{\dagger}_m(x)\ (1+i\alpha(x)\gamma^5)\
    \phi_n(x)\ a_n = \sum_n \bigl( \delta_{mn} + C_{mn} \bigr)a_n .
  \label{eq:19.68}
\end{equation}
Then
\begin{equation}
  \prod_m da'_m = \prod_m da_m\ \cdot\ J ,
  \label{eq:19.69}
\end{equation}
where $J$ is the Jacobian determinant of the transformation $(1+C)$.  The
inverse of $J$ appears in \eqref{eq:19.69} as a result of the rule
\eqref{eq:9.63} or Eq.~(9.69) for fermionic integration.  To evaluate
$J$, we write
\begin{equation}
  J = \det(1+C) = \exp\bigl[ \mathrm{tr}\log(1+C) \bigr]
  = \exp\Bigl\{ \sum_n C_{nn} + \cdots \Bigr\} ,
  \label{eq:19.70}
\end{equation}
and we can ignore higher order terms in the last line because $C$ is
infinitesimal.  Thus,
\begin{equation}
  \log J = i\int d^4x\ \alpha(x)\ \sum_n \phi^{\dagger}_n(x)\gamma^5\phi_n(x)
  .
  \label{eq:19.71}
\end{equation}
The coefficient of $\alpha(x)$ looks like $\mathrm{tr}[\gamma^5] = 0$.
However, we must regularize the sum over eigenstates $n$ in a gauge-invariant
way.  The natural choice is
\begin{equation}
  \sum_n \phi^{\dagger}_n(x)\gamma^5\phi_n(x)
  = \lim_{M\to\infty}\ \sum_n \phi^{\dagger}_n(x)\gamma^5\phi_n(x)\,
    e^{\lambda_n^2/M^2} .
  \label{eq:19.72}
\end{equation}
As \eqref{eq:19.65} indicates, the sign of $\lambda^2$ will be negative at
large momentum after a Wick rotation; thus, the sign in the exponent of the
convergence factor is given correctly.  We can write \eqref{eq:19.72} in an
operator form
\begin{equation}
  = \lim_{M\to\infty}\ \sum_n \phi^{\dagger}_n(x)\gamma^5\
    e^{(i\slashed{D})^2/M^2}\phi_n(x)
  = \lim_{M\to\infty}\ \mathrm{tr}\Bigl[ \gamma^5\ \langle x|\,
    e^{(i\slashed{D})^2/M^2}\,|x\rangle \Bigr] ,
  \label{eq:19.73}
\end{equation}
where, in the second line, we trace over Dirac indices.

To evaluate \eqref{eq:19.73}, we rewrite $(i\slashed{D})^2$ according to
\eqref{eq:16.107}.  In our present conventions, this equation reads
\begin{equation}
  (i\slashed{D})^2 = -D^2 + \frac{e}{2}\sigma^{\mu\nu}F_{\mu\nu} .
  \label{eq:19.74}
\end{equation}
Since we are taking the limit $M\to\infty$, we can concentrate our attention
on the asymptotic part of the spectrum, where the momentum $k$ is large and
we can expand in powers of the gauge field.  To obtain a nonzero trace with
$\gamma^5$, we must bring down four Dirac matrices from the exponent.  The
leading term is given by expanding the exponent to order $(\sigma\cdot F)^2$,
and then ignoring the background $A^\mu$ field in all other terms.  This gives
\begin{equation}
  \lim_{M\to\infty}\ \mathrm{tr}\Bigl[ \gamma^5\ \langle x|\,
    e^{(-D^2+(e/2)\sigma\cdot F)/M^2}\,|x\rangle \Bigr]
  = \lim_{M\to\infty}\ \mathrm{tr}\Bigl[ \gamma^5\ \frac{e^2}{8M^4}\,
    (\sigma\cdot F)^2\ \langle x|\,e^{-\partial^2/M^2}\,|x\rangle \Bigr] .
  \label{eq:19.75}
\end{equation}
The matrix element in \eqref{eq:19.75} can be evaluated by a Wick rotation:
\begin{equation}
  \langle x|\,e^{-\partial^2/M^2}\,|x\rangle
  = \int \frac{d^4k_E}{(2\pi)^4}\ e^{-k_E^2/M^2}
  = \frac{M^4}{16\pi^2} .
  \label{eq:19.76}
\end{equation}
Then \eqref{eq:19.75} reduces to
\begin{equation}
  \lim_{M\to\infty}\ \frac{-ie^2}{8\cdot16\pi^2}\ \mathrm{tr}\bigl[
    \gamma^5\sigma^{\mu\nu}\sigma^{\rho\sigma} \bigr]\ F_{\mu\nu}F_{\rho\sigma}
  = \frac{e^2}{32\pi^2}\ \epsilon^{\mu\nu\rho\sigma}F_{\mu\nu}F_{\rho\sigma}
  .
  \label{eq:19.77}
\end{equation}
Thus,
\begin{equation}
  J = \exp\Bigl[ i\int d^4x\ \alpha(x)\ \frac{e^2}{32\pi^2}\,
    \epsilon^{\mu\nu\rho\sigma}F_{\mu\nu}F_{\rho\sigma}(x) \Bigr] .
  \label{eq:19.78}
\end{equation}

In all, we find that, after the change of variables \eqref{eq:19.62}, the
functional integral \eqref{eq:19.61} takes the form
\begin{equation}
  Z = \int \mathcal{D}\bar{\psi}\,\mathcal{D}\psi\,\mathcal{D}A\ \exp\Bigl[
    i\int d^4x\ \Bigl( \mathcal{L} + \alpha(x)\Bigl( \partial_\mu j^{\mu 5}
      + \frac{e^2}{32\pi^2}\epsilon^{\mu\nu\rho\sigma}F_{\mu\nu}F_{\rho\sigma}
      \Bigr) \Bigr) \Bigr] .
  \label{eq:19.79}
\end{equation}
Varying the exponent with respect to $\alpha(x)$, we find precisely the
Adler-Bell-Jackiw anomaly equation.

This derivation of the axial vector anomaly is especially interesting because
it generalizes readily to any even dimensionality.  The functional derivation
always picks out for the right-hand side of the anomaly equation the
pseudoscalar operator built from gauge fields that has the same dimension,
$d$, as the divergence of the current.  In two dimensions, this derivation
leads immediately to \eqref{eq:19.18}.  As long as $d$ is even, we can always
construct a matrix $\gamma^5$ that anticommutes with all of the Dirac
matrices by taking their product.  Then, the functional derivation leads
straightforwardly to the result
\begin{equation}
  \partial_\mu j^{\mu 5} = \frac{1}{(2\pi)^n}\ \epsilon^{\mu_1\cdots\mu_{2n}}\
    F_{\mu_1\mu_2}\cdots F_{\mu_{2n-1}\mu_{2n}} ,
  \label{eq:19.80}
\end{equation}
where $n = d/2$.

At the end of the previous section, we argued that the axial vector anomaly
leads to global nonconservation of fermionic charges in a two-dimensional
system with a macroscopic electric field.  In the same way, the four-dimensional
anomaly equation leads to global nonconservation of the number of left- and
right-handed fermions in background fields in which the right-hand side of
\eqref{eq:19.45} is nonzero.  These are field configurations with parallel
electric and magnetic fields.  In Problem~19.1, we work out an example of
four-dimensional massless fermions in a simple situation of this type and
show that the fermion numbers are indeed violated, in a manner similar to
what we saw at the end of Section~19.1, in accord with the Adler-Bell-Jackiw
anomaly.

\section{Goldstone Bosons and Chiral Symmetries in QCD}

The Adler-Bell-Jackiw anomaly has a number of important implications for QCD.
To describe these, we must first discuss the chiral symmetries of QCD
systematically.  In this discussion, we will ignore all but the lightest
quarks $u$ and $d$.  In many analyses of the low-energy structure of the
strong interactions, one also treats the $s$ quark as light; this gives
results that naturally generalize the ones we will find below.

The fermionic part of the QCD Lagrangian is
\begin{equation}
  \mathcal{L} = \bar{u}i\slashed{D}u + \bar{d}i\slashed{D}d
    - m_u\bar{u}u - m_d\bar{d}d .
  \label{eq:19.81}
\end{equation}
If the $u$ and $d$ quarks are very light, the last two terms are small and
can be neglected.  Let us study the implications of making this
approximation.  If we ignore the $u$ and $d$ masses, the Lagrangian
\eqref{eq:19.81} of course has isospin symmetry, the symmetry of an $SU(2)$
unitary transformation mixing the $u$ and $d$ fields.  However, because the
classical Lagrangian for massless fermions contains no coupling between left-
and right-handed quarks, this Lagrangian actually is symmetric under the
separate unitary transformations
\begin{equation}
  q_L \to U_L q_L,\qquad q_R \to U_R q_R .
  \label{eq:19.82}
\end{equation}
It is useful to separate the $U(1)$ and $SU(2)$ parts of these
transformations; then the symmetry group of the classical, massless QCD
Lagrangian is $SU(2)\times SU(2)\times U(1)\times U(1)$.  Let $q$ denote the
quark doublet, with chiral components
\begin{equation}
  q = \begin{pmatrix}u\\ d\end{pmatrix},\qquad
  q_L = \frac{1-\gamma^5}{2}q,\qquad q_R = \frac{1+\gamma^5}{2}q .
  \label{eq:19.83}
\end{equation}
Then we can write the currents associated with these symmetries as
\begin{equation}
  j^\mu_L = \bar{q}_L\gamma^\mu q_L,\qquad j^\mu_R = \bar{q}_R\gamma^\mu q_R,
  \qquad j^{\mu a}_L = \bar{q}_L\gamma^\mu\tau^a q_L,\qquad
  j^{\mu a}_R = \bar{q}_R\gamma^\mu\tau^a q_R ,
  \label{eq:19.84}
\end{equation}
where $\tau^a = \sigma^a/2$ represent the generators of $SU(2)$.  The sums of
left- and right-handed currents give the baryon number and isospin currents
\begin{equation}
  j^\mu = j^\mu_L + j^\mu_R,\qquad j^{\mu a} = j^{\mu a}_L + j^{\mu a}_R .
  \label{eq:19.85}
\end{equation}
The corresponding symmetries are the transformations \eqref{eq:19.82} with
$U_L = U_R$.  The differences of the currents \eqref{eq:19.84} give the
corresponding axial vector currents $j^{\mu a}_5$:
\begin{equation}
  j^{\mu a}_5 = j^{\mu a}_R - j^{\mu a}_L .
  \label{eq:19.86}
\end{equation}
In the discussion to follow, we will derive conclusions about the strong
interactions by assuming that the classical conservation laws for these
currents are not spoiled by anomalies.  We will show below that this
assumption is correct for the isotriplet currents $j^{\mu a}_5$ but not for
$j^{\mu 5}$.

The vector $SU(2)\times U(1)$ transformations are manifest symmetries of the
strong interactions.  The axial transformations, however, are not symmetries
of the ground state: they are spontaneously broken.  The vacuum expectation
value of the operator $\bar{q}q$ is nonzero,
\begin{equation}
  \langle 0|\bar{u}u|0\rangle = \langle 0|\bar{d}d|0\rangle \neq 0 ,
  \label{eq:19.87}
\end{equation}
which breaks the axial symmetry.  The Goldstone bosons associated with the
spontaneous breaking of the $SU(2)$ axial symmetry are the pions.  The
current $j^{\mu a}_5$ creates a pion from the vacuum:
\begin{equation}
  \langle 0|j^{\mu a}_5(x)|\pi^b(p)\rangle = -ip^\mu f_\pi\,\delta^{ab}\,
    e^{-ip\cdot x} ,
  \label{eq:19.88}
\end{equation}
where $f_\pi$ is the pion decay constant.

The spontaneous breaking of chiral symmetry has many consequences.  In
particular, it implies relations among the pion mass, the quark masses, and
the matrix elements of the axial currents.  The pion is a Goldstone boson of
a broken symmetry, and its interactions are determined by the low-energy
theorems that follow from this fact.  To illustrate these ideas, let us
consider the axial current vertex between nucleon states.  We write the
nucleon matrix element of the axial isospin current as
\begin{equation}
  \langle N(p')|j^{\mu a}_5|N(p)\rangle = \bar{u}(p')\ \Bigl[
    \gamma^\mu\gamma^5 g_A(q^2) + \frac{q^\mu}{2m_N}\gamma^5\ g_P(q^2)
    \Bigr]\ \tau^a\,u(p) ,
  \label{eq:19.97}
\end{equation}
where $q = p'-p$.  The conservation of the axial current implies
\begin{equation}
  q_\mu\langle N(p')|j^{\mu a}_5|N(p)\rangle = 0 .
  \label{eq:19.98}
\end{equation}
This equation implies that $g_A$ must vanish unless the current matrix
element contains a pole in $q^2$.  Such a pole would imply the presence of a
physical massless particle, but fortunately, there is one available---the
massless pion.  The process in which the current creates a pion that is then
absorbed by the nucleon indeed leads to a pole in the form factor, as shown
in Fig.~\ref{fig:19.6}(b).

\begin{figure}[t]
  \centering
  \includegraphics[width=0.55\textwidth]{fig19.6.png}
  \caption{Diagrams contributing to the nucleon matrix element of the axial
    current: (a) the direct coupling; (b) the pion pole contribution.}
  \label{fig:19.6}
\end{figure}

Let us now compute this pole term and use it to determine $g_A$.  The
low-energy pion-nucleon interaction is conventionally parametrized by the
Lagrangian
\begin{equation}
  \Delta\mathcal{L} = ig_{\pi NN}\,\bar{N}\gamma^5\tau^a N\,\pi^a .
  \label{eq:19.99}
\end{equation}
The amplitude for the current $j^{\mu a}_5$ to create the pion is given by
\eqref{eq:19.88}.  Then the contribution of Fig.~\ref{fig:19.6}(b) to the
current vertex is
\begin{equation}
  \frac{1}{q^2}\ (if_\pi q^\mu)\ g_{\pi NN}\ \bar{N}\gamma^5\tau^a N .
  \label{eq:19.100}
\end{equation}
Thus,
\begin{equation}
  g_P(q^2) = \frac{2m_N f_\pi g_{\pi NN}}{q^2} .
  \label{eq:19.101}
\end{equation}
Using the conservation condition \eqref{eq:19.98}, we find that $g_A$ is
given by a combination of the pion decay constant, the nucleon mass, and the
pion-nucleon coupling constant:
\begin{equation}
  g_A = \frac{f_\pi g_{\pi NN}}{m_N} .
  \label{eq:19.102}
\end{equation}
This strange identity, called the \emph{Goldberger-Treiman relation}, is
satisfied experimentally to 5\% accuracy.

The identification of the pion as the Goldstone boson of spontaneously
broken chiral symmetry leads to numerous other predictions for current
matrix elements and pion scattering amplitudes.  In particular, the leading
terms of the pion-pion and pion-nucleon scattering amplitudes at low energy
can be computed directly in terms of $f_\pi$ by arguments similar to the one
just given.%
\footnote{The detailed consequences of spontaneously broken chiral symmetry
are worked out in a very clear manner in Georgi (1984).}

\subsection{Anomalies of Chiral Currents}

Up to this point, we have discussed the chiral symmetries of QCD according to
the classical current conservation equations.  We must now ask whether these
equations are affected by the Adler-Bell-Jackiw anomaly, and what the
consequences of that modification are.

\begin{figure}[t]
  \centering
  \includegraphics[width=0.5\textwidth]{fig19.7.png}
  \caption{Diagrams that lead to an axial vector anomaly for a chiral
    current in QCD.}
  \label{fig:19.7}
\end{figure}

To begin, we study the modification of the chiral conservation laws due to
the coupling of the quark currents to the gluon fields of QCD.  The arguments
given in the previous section go through equally well in the case of massless
fermions coupling to a non-Abelian gauge field, so we expect that an axial
vector current will receive an anomalous contribution from the diagrams shown
in Fig.~\ref{fig:19.7}.  The anomaly equation should be the Abelian result,
supplemented by an appropriate group theory factor.  In addition, since the
axial current is gauge invariant, the anomaly must also be gauge invariant.
That is, it must contain the full non-Abelian field strength, including its
nonlinear terms.  These terms are actually included in the functional
derivation of the anomaly given at the end of Section~19.2.

For the axial currents of QCD, written in \eqref{eq:19.86}, we can read the
group theory factors for the Adler-Bell-Jackiw anomaly from the diagrams of
Fig.~\ref{fig:19.7}.  For the axial isospin currents,
\begin{equation}
  \partial_\mu j^{\mu a}_5 = \frac{g^2}{32\pi^2}\ \epsilon^{\mu\nu\rho\sigma}\
    F^c_{\mu\nu}F^d_{\rho\sigma}\ \mathrm{tr}\bigl[ \tau^a\{t^c, t^d\} \bigr],
  \label{eq:19.103}
\end{equation}
where $F^c_{\mu\nu}$ is a gluon field strength, $\tau^a$ is an isospin
matrix, $t^c$ is a color matrix, and the trace is taken over colors and
flavors.  In this case, we find
\begin{equation}
  \mathrm{tr}\bigl[ \tau^a t^c t^d \bigr] = \mathrm{tr}[\tau^a]\,
    \mathrm{tr}[t^c t^d] = 0 ,
  \label{eq:19.104}
\end{equation}
since the trace of a single $\tau^a$ vanishes.  Thus the conservation of the
axial isospin currents is unaffected by the Adler-Bell-Jackiw anomaly of QCD.
However, in the case of the isospin singlet axial current, the matrix $\tau^a$
is replaced by the matrix 1 on flavors, and we find
\begin{equation}
  \partial_\mu j^{\mu 5} = n_f\ \frac{g^2}{32\pi^2}\ \epsilon^{\mu\nu\rho\sigma}
    F^c_{\mu\nu}F^c_{\rho\sigma} ,
  \label{eq:19.105}
\end{equation}
where $n_f$ is the number of flavors; $n_f = 2$ in our current model.

Thus, the isospin singlet axial current is not in fact conserved in QCD.  The
divergence of this current is equal to a gluon operator with nontrivial matrix
elements between hadron states.  Some subtle questions remain concerning the
effects of this operator.  In particular, it can be shown, as we saw for the
two-dimensional axial anomaly in Eq.~(19.31), that the right-hand
side of \eqref{eq:19.105} is a total divergence.  Nevertheless, again in
accord with our experience in two dimensions, there are physically reasonable
field configurations in which the four-dimensional integral of this term
takes a nonzero value.  This topic is discussed further at the end of
Section~22.3.  In any event, Eq.~\eqref{eq:19.105} indeed implies that QCD
has no isosinglet axial symmetry and no associated Goldstone boson.  This
equation explains why the strong interactions contain no light isosinglet
pseudoscalar meson with mass comparable to that of the pions.

Though the axial isospin currents have no axial anomaly from QCD
interactions, they do have an anomaly associated with the coupling of quarks
to electromagnetism.  Again referring to the diagrams of Fig.~\ref{fig:19.7},
we see that the electromagnetic anomaly of the axial isospin currents is given
by
\begin{equation}
  \partial_\mu j^{\mu a}_5 = \frac{e^2}{32\pi^2}\ \epsilon^{\mu\nu\rho\sigma}\
    F_{\mu\nu}F_{\rho\sigma}\ \mathrm{tr}\bigl[ \tau^a Q^2 \bigr] ,
  \label{eq:19.106}
\end{equation}
where $F_{\mu\nu}$ is the electromagnetic field strength, $Q$ is the matrix
of quark electric charges,
\begin{equation}
  Q = \begin{pmatrix}\frac{2}{3}&0\\ 0&-\frac{1}{3}\end{pmatrix},
  \label{eq:19.107}
\end{equation}
and the trace again runs over flavors and colors.  Since the matrices in the
trace do not depend on color, the color sum simply gives a factor of 3.  The
flavor trace is nonzero only for $a = 3$; in that case, the electromagnetic
anomaly is
\begin{equation}
  \partial_\mu j^{\mu 3}_5 = \frac{e^2}{16\pi^2}\ \epsilon^{\mu\nu\rho\sigma}\
    F_{\mu\nu}F_{\rho\sigma} .
  \label{eq:19.108}
\end{equation}
Because the current $j^{\mu 3}_5$ annihilates a $\pi^0$ meson, Eq.~\eqref{eq:19.108}
indicates that the axial vector anomaly contributes to the matrix element for
the decay $\pi^0\to 2\gamma$.  We will now show that, in fact, it gives the
leading contribution to this amplitude.

Again, we work in the limit of massless $u$ and $d$ quarks, so that the
chiral symmetries are exact up to the effects of the anomaly.  Consider the
matrix element of the axial current between the vacuum and a two-photon
state:
\begin{equation}
  \langle \gamma(p)\gamma(k)|j^{\mu 5}(q)|0\rangle = M^{\mu\nu\lambda}
    \epsilon_\nu(p)\epsilon_\lambda(k) .
  \label{eq:19.109}
\end{equation}
This is the same matrix element Eq.~(19.46) that we studied in QED
perturbation theory in Section~19.2.  Now, however, we will study the general
properties of this matrix element by expanding it in form factors.  In
general, the amplitude can be decomposed by writing all possible tensor
structures and applying the restrictions that follow from symmetry under the
interchange of $(\mu,\nu)$ and $(k,\lambda)$ and the QED Ward identities
Eq.~(19.51).  This leaves three possible structures:
\begin{equation}
  \begin{aligned}
    M^{\mu\nu\lambda} &= \epsilon^{\mu\nu\lambda\alpha} p_\alpha k^\nu\,
      M_1 + \bigl( \epsilon^{\mu\nu\alpha\beta}k^\lambda - \epsilon^{\mu\lambda
        \alpha\beta}p^\nu \bigr)\,k_\alpha p_\beta\,M_2\\
    &\quad + \Bigl[ \bigl( \epsilon^{\mu\nu\alpha\beta}p^\lambda
      - \epsilon^{\mu\lambda\alpha\beta}k^\nu \bigr)\,k_\alpha p_\beta
      \cdot k \Bigr]\,M_3 .
  \end{aligned}
  \label{eq:19.110}
\end{equation}
The second term satisfies Eq.~(19.51) by virtue of the on-shell
conditions $p^2 = k^2 = 0$.

Now contract \eqref{eq:19.110} with $(iq_\mu)$ to take the divergence of the
axial vector current.  We find
\begin{equation}
  iq_\mu M^{\mu\nu\lambda} = iq^2\ \epsilon^{\nu\lambda\alpha\beta}\ p_\alpha
    k_\beta\ M_1 - i\ \epsilon^{\nu\lambda\alpha\beta}\ q_\alpha(p-k)_\beta\,
    k\cdot p\ M_3 ;
  \label{eq:19.111}
\end{equation}
the other terms automatically give zero.  Using $q = p + k$ and
$q^2 = 2p\cdot k$, we can simplify this to
\begin{equation}
  iq_\mu M^{\mu\nu\lambda} = iq^2\ \epsilon^{\nu\lambda\alpha\beta}\ p_\alpha
    k_\beta\ \bigl( M_1 + M_3 \bigr) .
  \label{eq:19.112}
\end{equation}
The whole quantity is proportional to $q^2$ and apparently vanishes in the
limit $q^2\to 0$.  This contrasts with the prediction of the axial vector
anomaly.  Taking the matrix element of the right-hand side of
\eqref{eq:19.108}, we find
\begin{equation}
  iq_\mu M^{\mu\nu\lambda} = \frac{e^2}{4\pi^2}\ \epsilon^{\nu\lambda\alpha\beta}
    p_\alpha k_\beta .
  \label{eq:19.113}
\end{equation}
The conflict can be resolved if one of the form factors appearing in
\eqref{eq:19.112} contains a pole in $q^2$.  Such a pole can arise through
the process shown in Fig.~\ref{fig:19.8}, in which the current creates a
$\pi^0$ meson which subsequently decays to two photons.  The amplitude for
the current to create the meson is given by \eqref{eq:19.88}.  Let us
parametrize the pion decay amplitude as
\begin{equation}
  i\mathcal{M}(\pi^0 \to 2\gamma) = iA\,e\ \epsilon^{\nu\lambda\alpha\beta}\,
    p_\alpha k_\beta\ \epsilon_\nu(p)\epsilon_\lambda(k) ,
  \label{eq:19.114}
\end{equation}
where $A$ is a constant to be determined.  Then the contribution of the
process of Fig.~\ref{fig:19.8} to the amplitude $M^{\mu\nu\lambda}$ defined in
\eqref{eq:19.109} is
\begin{equation}
  \bigl( iq^\mu f_\pi \bigr)\ \frac{1}{q^2}\ \bigl( iA\,e\,
    \epsilon^{\nu\lambda\alpha\beta}p_\alpha k_\beta \bigr) .
  \label{eq:19.115}
\end{equation}
This is a contribution to the form factor $M_1$,
\begin{equation}
  M_1 = \frac{f_\pi}{q^2}\cdot A\,e ,
  \label{eq:19.116}
\end{equation}
plus terms regular at $q^2 = 0$.  Now, by equating \eqref{eq:19.112} to
\eqref{eq:19.113}, we determine $A$ in terms of the coefficient of the
anomaly:
\begin{equation}
  A = \frac{1}{4\pi^2 f_\pi} .
  \label{eq:19.117}
\end{equation}
From the decay matrix element \eqref{eq:19.114}, it is straightforward to
work out the decay rate of the $\pi^0$.  Note that, though we have worked out
the decay matrix element in the limit of a massless $\pi^0$, we must supply
the physically correct kinematics which depends on the $\pi^0$ mass.
Including a factor 1/2 for the phase space of identical particles, we find
\begin{equation}
  \Gamma(\pi^0 \to 2\gamma) = \frac{1}{2}\ \frac{1}{32\pi m_\pi}\ \int
    d\Omega\ |\mathcal{M}|^2 = \frac{\alpha^2 m_\pi^3}{64\pi^3 f_\pi^2} .
  \label{eq:19.118}
\end{equation}
This formula is in excellent agreement with experiment.  The prediction of
the $\pi^0$ lifetime from the axial anomaly, combined with the measured
value of $f_\pi$, is a striking confirmation of the Adler-Bell-Jackiw
anomaly.

\begin{figure}[t]
  \centering
  \includegraphics[width=0.5\textwidth]{fig19.8.png}
  \caption{Contribution that leads to a pole in the axial vector current
    form factor $M_1$.}
  \label{fig:19.8}
\end{figure}

\section{Chiral Anomalies and Chiral Gauge Theories}

The anomalies that we have discussed so far arise in theories with vector-like
couplings, such as QED and QCD.  In these theories, the anomaly affects only
the axial currents, which are not related to the gauge symmetry, and so the
consistency of the theory is not endangered.  In a \emph{chiral gauge
theory}, however, the gauge fields couple differently to left- and
right-handed fermions.  In such theories, the axial anomaly can affect the
conservation of the gauge currents themselves, which would destroy the
consistency of the theory.  We will now study this problem.

Consider a chiral gauge theory with gauge group $G$, in which the
left-handed fermions transform in a representation $R$ of $G$.  A typical
example is a theory of left-handed quarks and leptons, with the weak
interaction gauge group $SU(2)$ coupled only to left-handed doublets:
\begin{equation}
  \psi_L = \begin{pmatrix}\nu_L\\ e_L\end{pmatrix},\qquad
  q_L = \begin{pmatrix}u_L\\ d_L\end{pmatrix},
  \label{eq:19.124}
\end{equation}
and no corresponding right-handed couplings.

The gauge currents of such a theory receive corrections from the triangle
diagrams shown in Fig.~\ref{fig:19.9}.  These diagrams are the analogues of
the anomaly diagrams that we studied in Section~19.2, but now the external
currents are gauge currents.  The anomaly in the divergence of the gauge
current $j^{\mu a}$ is given by
\begin{equation}
  \partial_\mu j^{\mu a} = \frac{g^2}{32\pi^2}\ \epsilon^{\mu\nu\rho\sigma}\
    F^b_{\mu\nu}F^c_{\rho\sigma}\ A^{abc} ,
  \label{eq:19.131}
\end{equation}
where
\begin{equation}
  A^{abc} = \mathrm{tr}\bigl[ t^a\{t^b, t^c\} \bigr] ,
  \label{eq:19.132}
\end{equation}
with the trace taken over the fermion representation $R$.

\begin{figure}[t]
  \centering
  \includegraphics[width=0.5\textwidth]{fig19.9.png}
  \caption{Diagrams contributing to the anomaly of a gauge symmetry current
    in a chiral gauge theory.}
  \label{fig:19.9}
\end{figure}

This equation implies that, unless $A^{abc}$ vanishes, the current
$j^{\mu a}$ is not conserved.  The factor \eqref{eq:19.132} is totally
symmetric in $(a, b, c)$, so this condition is independent of which current
is treated as an external operator.  As we described in Sections~19.1
and~19.2, we can change the regularization of the diagram so that the
external current is conserved, but only at the price of violating the
conservation of one of the other two currents in the diagram.

Since the whole construction of a theory with local gauge invariance is based
on the existence of an exact global symmetry, the violation of the
conservation of $j^{\mu a}$ does violence to the structure of the theory.
For example, triangle diagrams of the form of Fig.~\ref{fig:19.9} will now
generate divergent gauge boson mass terms and will upset the delicate
relations between three- and four-point vertices discussed in Chapter~16.
These relations, following from the Ward identity, were necessary to insure
the cancellation of unphysical states and the unitarity of the $S$-matrix.
The only way to avoid this problem is to insist that $A^{abc} = 0$ as a
fundamental consistency condition for chiral gauge theories.%
\footnote{D.~J.~Gross and R.~Jackiw, \emph{Phys.\ Rev.\ D} \textbf{6}, 477
(1972).}  Gauge theories satisfying this condition are said to be
\emph{anomaly free}.

As an example of the application of this condition, consider the prototype
weak interaction gauge theory that we presented in \eqref{eq:19.124}.  If the
two gauge bosons in Fig.~\ref{fig:19.9} are $SU(2)$ gauge bosons and the
current $j^{\mu a}$ is an $SU(2)$ gauge current, we would evaluate
\eqref{eq:19.132} by substituting $t^a = \tau^a = \sigma^a/2$ and using the
relation $\{\sigma^b, \sigma^c\} = 2\delta^{bc}$.  This gives
\begin{equation}
  A^{abc} = \mathrm{tr}\bigl[ \sigma^a\cdot 2\delta^{bc} \bigr] = 0 ,
  \label{eq:19.133}
\end{equation}
so the consistency condition is satisfied.  If the fermions in
\eqref{eq:19.124} also couple to electromagnetism, there is an additional
consistency condition that we would find by taking the current in
Fig.~\ref{fig:19.9} to be the electromagnetic current.  The factor $A^{abc}$
for this case is
\begin{equation}
  \mathrm{tr}\bigl[ Q\{t^b, t^c\} \bigr] ,
  \label{eq:19.134}
\end{equation}
where $Q$ is the matrix of electric charges.  If we simplify as in
\eqref{eq:19.133}, the trace \eqref{eq:19.134} becomes
\begin{equation}
  \frac{1}{2}\mathrm{tr}[Q]\,\delta^{bc} .
  \label{eq:19.135}
\end{equation}
This factor is proportional to the sum of the fermion electric charges, which
does not vanish either for quarks or for leptons.  However, if we sum over one
quark doublet and one lepton doublet, with a factor 3 for colors, we find
\begin{equation}
  \mathrm{tr}[Q] = 3\cdot\Bigl( \frac{2}{3} - \frac{1}{3} \Bigr)
    + (0 - 1) = 0 .
  \label{eq:19.136}
\end{equation}
Remarkably, the weak interaction gauge theory described by \eqref{eq:19.124}
can be consistently combined with QED only if the theory contains equal
numbers of quark and lepton doublets.

We complete this section by working out more generally the condition that a
chiral gauge theory be anomaly free.  We will first derive some basic
properties of the anomaly factor $A^{abc}$ and then apply these to chiral
gauge theories with simple gauge groups.

If the fermion representation $R$ is real, $R$ is equivalent to its conjugate
representation $\bar{R}$.  Thus, as we described below \eqref{eq:15.82},
$t^a_R$ is related by a unitary transformation to
$t^a_{\bar{R}} = -(t^a_R)^T$.  Since \eqref{eq:19.132} is invariant to
unitary transformations of the $t^a$, we can replace $t^a_R$ by $t^a_{\bar{R}}$.
Then
\begin{equation}
  \begin{aligned}
    A^{abc} &= \mathrm{tr}\Bigl[ (-t^a)^T\ \bigl\{ (-t^b)^T, (-t^c)^T
      \bigr\} \Bigr]\\
    &= -\mathrm{tr}\Bigl[ \{ t^c, t^b \}\ t^a \Bigr] = -A^{abc} .
  \end{aligned}
  \label{eq:19.137}
\end{equation}
Thus, if $R$ is real, the gauge theory is automatically anomaly free.  As a
special case, any gauge theory of Dirac fermions is anomaly free.

In more general circumstances, we can simplify the calculation of $A^{abc}$
by noting that it is an invariant of the gauge group $G$ that is totally
symmetric with three indices in the adjoint representation.  For some
possible groups, a suitable invariant may not exist, and in those cases
$A^{abc}$ must vanish.  For example, in $SU(2)$ the adjoint representation has
spin 1.  The symmetric product of two spin-1 multiplets gives spin 0 plus
spin 2, with no spin-1 component.  Thus, there is no symmetric tensor
coupling two spin-1 indices to give a spin 1.  The factor \eqref{eq:19.132}
must then vanish in any $SU(2)$ gauge theory.  We saw this happen in an
explicit example in Eq.~\eqref{eq:19.133}.

In $SU(n)$ groups, $n > 3$, there is a unique symmetric invariant $d^{abc}$
of the required type.  It appears in the anticommutator of representation
matrices of the fundamental representation:
\begin{equation}
  \{ t^a_n, t^b_n \} = \frac{1}{n}\delta^{ab} + d^{abc}t^c_n .
  \label{eq:19.138}
\end{equation}
The uniqueness of this invariant implies that, in an $SU(n)$ gauge theory,
any trace of the form of \eqref{eq:19.132} is proportional to $d^{abc}$.
For each representation $r$, we can define an anomaly coefficient $A(r)$ by
\begin{equation}
  \mathrm{tr}\bigl[ t^a_r\{t^b_r, t^c_r\} \bigr] = \frac{1}{2}A(r)\,d^{abc}
  .
  \label{eq:19.139}
\end{equation}
For the fundamental representation, we can see from \eqref{eq:19.138} that
$A(n) = 1$.  It follows from the argument of \eqref{eq:19.137} that
\begin{equation}
  A(\bar{r}) = -A(r) .
  \label{eq:19.140}
\end{equation}
For higher representations, the anomaly coefficients can be worked out using
methods similar to those we used in Section~15.4 to compute $C_2(r)$.  For
example, we show in Problem~19.3 that, if $a$ and $s$ are the $SU(n)$
representations corresponding to antisymmetric and symmetric two-index
tensors, then
\begin{equation}
  A(a) = n - 4,\qquad A(s) = n + 4 .
  \label{eq:19.141}
\end{equation}
An $SU(n)$ gauge theory is anomaly free if the anomaly coefficients of the
various irreducible components of the fermion multiplet $R$ sum to zero.  For
example, the $SU(n)$ gauge theory of left-handed fermions with
representation content
\begin{equation}
  R = a + (n-4)\bar{n}
  \label{eq:19.142}
\end{equation}
is anomaly free.

Of the various simple Lie groups listed below \eqref{eq:15.72}, only
$SU(n)$, $SO(4n+2)$, and $E_6$ have complex representations.  Of these, only
$SU(n)$ and $SO(6)$, which has the same Lie algebra as $SU(4)$, have a
symmetric invariant of the type required to build the anomaly.  Gauge
theories based on $SO(4n+2)$, $n > 2$, and on $E_6$ are automatically
anomaly free.  The groups $SO(10)$ and $E_6$ have been suggested as
candidates for the grand unified gauge symmetry of particle physics, which we
will discuss in Section~22.2.

There is one further constraint on the representation content of a chiral
gauge theory, which comes from considering its coupling to gravity.  It is
possible to show that the diagrams of Fig.~\ref{fig:19.9} give an anomaly
contribution when computed with a gauge current $j^{\mu a}$ and external
gravitational fields.  The group-theory factor that multiplies this diagram
is
\begin{equation}
  \mathrm{tr}\bigl[ t^a \bigr] .
  \label{eq:19.143}
\end{equation}
This factor automatically vanishes if the gauge group is non-Abelian.
However, if the gauge group of the theory contains $U(1)$ factors, the theory
cannot be consistently coupled to gravity unless each of the $U(1)$
generators is traceless.%
\footnote{L.~Alvarez-Gaum\'e and E.~Witten, \emph{Nucl.\ Phys.\ B} \textbf{234},
269 (1984).}

Once we have constructed a consistent chiral gauge theory, we have an
additional problem of finding a prescription for calculating in this theory
consistently.  In a vector-like gauge theory, we can define ultraviolet-divergent
diagrams with dimensional regularization.  This guarantees that the divergent
diagrams will be regulated in a way that respects the Ward identities of
local gauge invariance.  To generalize dimensional regularization to chiral
gauge theories, we need to introduce a dimensional continuation of
$\gamma^5$.  The 't Hooft-Veltman definition used to define the chiral current
in Section~19.2 is not satisfactory, because this definition does not
manifestly respect the conservation of the gauge currents.  A useful
alternative procedure is to define $\gamma^5$ formally as an object that
anticommutes with all of the $\gamma^\mu$.  This prescription gives
unambiguous, gauge-invariant results for amplitudes that are not proportional
to $\epsilon^{\mu\nu\rho\sigma}$, at least through two-loop order.  In
Section~21.3, we will use this prescription to compute loop diagrams in weak
interaction theory.  As a last resort, one can always compute with a
non-gauge-invariant regulator and add non-gauge-invariant counterterms to the
theory so that the gauge theory Ward identities remain valid.

\section{Anomalous Breaking of Scale Invariance}

There is one more important example of a symmetry that is an invariance at
the classical level and is broken by quantum corrections.  This is the
classical \emph{scale invariance} of a massless field theory with a
dimensionless coupling constant.  In Chapter~12, we saw that a quantum field
theory with no classical dimensionful parameters still depends on a mass
scale through the regularization of ultraviolet divergences, or, equivalently,
through the running of coupling constants.  We have already seen how to
analyze this induced dependence on the renormalization scale using the
Callan-Symanzik equation.  In this section, we will show how the violation
of classical scale invariance by quantum corrections can be described as a
current conservation anomaly.

In this book we have avoided giving a careful treatment of the energy-momentum
tensor of a quantum field theory.  In Section~2.2, we used Noether's theorem
to demonstrate that the invariance of a quantum field theory under spacetime
translations implies the presence of a conserved tensor $T^{\mu\nu}$.  In
Section~9.6, we gave an alternative derivation of this result using the
functional integral formalism.  However, to discuss the theory of scale
invariance, we will need some more detailed properties of the energy-momentum
tensor.  We will now simply state these properties and refer elsewhere for
their derivations.%
\footnote{The conclusions presented in the next three paragraphs are derived
with care in C.~G.~Callan, S.~Coleman, and R.~Jackiw, \emph{Ann.\ Phys.}
\textbf{59}, 42 (1970).}

The tensor $T^{\mu\nu}$ defined by expressions \eqref{eq:2.17} and
\eqref{eq:9.99} is called the \emph{canonical energy-momentum tensor}.  The
expressions that defined this tensor do not imply that $T^{\mu\nu}$ is
symmetric.  In fact, this tensor need not be symmetric, and, in a gauge
theory, it need not be gauge-invariant.  However, it is always possible to
convert $T^{\mu\nu}$ into a symmetric and gauge-invariant tensor $\Theta^{\mu\nu}$
by the addition
\begin{equation}
  \Theta^{\mu\nu} = T^{\mu\nu} + \partial_\lambda K^{\lambda\mu\nu} ,
  \label{eq:19.144}
\end{equation}
where $K^{\lambda\mu\nu}$ is antisymmetric under interchange of $\mu$ and
$\lambda$.  The form of the added term implies that $\Theta^{\mu\nu}$ is
conserved if $T^{\mu\nu}$ is, and that the global energy-momentum four-vector
is unchanged.

A scale transformation of a scalar field theory can be defined as a
transformation of variables
\begin{equation}
  \phi(x) \to e^{-D\alpha}\phi(x e^{-\alpha}) ,
  \label{eq:19.146}
\end{equation}
with $D = 1$, the canonical mass dimension of the field.  The scale
transformation is defined similarly in theories of fermion and gauge fields.
If this transformation is an invariance of the classical Lagrangian, as it
will be if there are no dimensionful couplings, this theory will possess a
conserved current $D^\mu$ called the \emph{dilatation current}.  The
dilatation current has a simple relation to the symmetric energy-momentum
tensor $\Theta^{\mu\nu}$: $D^\mu = x_\nu\Theta^{\mu\nu}$, so that
\begin{equation}
  \partial_\mu D^\mu = \Theta^\mu_{\ \mu} .
  \label{eq:19.147}
\end{equation}
The derivation of these results from Noether's theorem is not
straightforward.  There is a simpler derivation, which, however, uses
formalism beyond the scope of this book.  If the quantum field theory under
consideration is coupled to gravity, then the energy-momentum tensor can be
identified as the source of the gravitational field.  This energy-momentum
tensor can be found by varying the Lagrangian $\mathcal{L}$ of matter fields
with respect to the spacetime metric $g_{\mu\nu}(x)$.  This construction
gives a manifestly symmetric and gauge-invariant tensor, which turns out to
be
\begin{equation}
  \Theta^{\mu\nu} = \frac{2}{\sqrt{-g}}\ \frac{\delta\mathcal{L}}
    {\delta g_{\mu\nu}} .
  \label{eq:19.148}
\end{equation}
A scale transformation can be represented as a change in the spacetime metric
\begin{equation}
  g_{\mu\nu}(x) \to e^{2\alpha}g_{\mu\nu}(x) .
  \label{eq:19.149}
\end{equation}
Combining \eqref{eq:19.148} and \eqref{eq:19.149}, we see that the change in
the Lagrangian induced by this transformation is just the trace of
$\Theta^{\mu\nu}$.  This will be equal by Noether's theorem to the divergence
of the corresponding current, giving us back Eq.~\eqref{eq:19.147}.

In QED, it is not hard to guess the form of the symmetric energy-momentum
tensor:
\begin{equation}
  \Theta^{\mu\nu} = -F^{\mu\lambda}F^{\nu}_{\ \lambda}
    + \frac{1}{4}g^{\mu\nu}(F_{\lambda\sigma})^2
    + \bar{\psi}\gamma^{(\mu}iD^{\nu)}\psi
    - g^{\mu\nu}\bar{\psi}(i\slashed{D} - m)\psi .
  \label{eq:19.150}
\end{equation}
This is a gauge-invariant symmetric tensor that leads to the familiar
expression for the total energy,
\begin{equation}
  H = \int d^3x\ \Bigl\{ \frac{1}{2}(\mathbf{E}^2 + \mathbf{B}^2)
    + \psi^{\dagger}\bigl( -i\bm{\alpha}\cdot\nabla + m\beta \bigr)\psi
    \Bigr\} .
  \label{eq:19.151}
\end{equation}
For future reference, we note that these results are true at the classical
level in any spacetime dimension $d$.  In four dimensions, the trace of the
gauge field terms cancels automatically.  After using the Dirac equation,
which is valid as an operator equation of motion, we find that the trace of
$\Theta^{\mu\nu}$ is given by
\begin{equation}
  \Theta^\mu_{\ \mu} = m\bar{\psi}\psi ,
  \label{eq:19.152}
\end{equation}
and indeed vanishes, classically, if $m = 0$.

When quantum corrections are included, we know that a scale transformation is
not a symmetry of the theory, since the same theory referred to a larger
scale contains a different value of the renormalized coupling constant.  The
shift of the renormalized coupling is
\begin{equation}
  g \to g + \alpha\beta(g) ,
  \label{eq:19.153}
\end{equation}
and the corresponding change in the Lagrangian is
\begin{equation}
  \delta\mathcal{L} = \alpha\beta(g)\,\frac{\partial\mathcal{L}}
    {\partial g} .
  \label{eq:19.154}
\end{equation}
Thus, when quantum corrections are included, the equation for the dilatation
current in a classically scale-invariant theory should read
\begin{equation}
  \partial_\mu D^\mu = \Theta^\mu_{\ \mu} = \beta(g)\,\frac{\partial
    \mathcal{L}}{\partial g} .
  \label{eq:19.155}
\end{equation}
In massless QED, we can write this formula most usefully by rescaling the
gauge fields so that the coupling constant $e$ is removed from the covariant
derivative: $eA_\mu\to A_\mu$.  Then $e$ appears only in the term
\begin{equation}
  \mathcal{L} \supset -\frac{1}{4e^2}(F_{\lambda\sigma})^2 + \cdots ,
  \label{eq:19.156}
\end{equation}
and Eq.~\eqref{eq:19.155} reads
\begin{equation}
  \Theta^\mu_{\ \mu} = \frac{\beta(e)}{2e}\ \frac{1}{2}(F_{\lambda\sigma})^2
  = \frac{e^2}{24\pi^2}\ \frac{1}{4}(F_{\lambda\sigma})^2 .
  \label{eq:19.157}
\end{equation}
This relation, which says that the trace of the symmetric energy-momentum
tensor takes a nonzero value as a result of quantum corrections, is known as
the \emph{trace anomaly}.

We should be able to check the trace anomaly equation \eqref{eq:19.157}
directly in a computation.  The physical content of the trace anomaly is that
the trace of the energy-momentum tensor, which vanishes classically, receives
a nonzero contribution from the running of the coupling constant.  This
anomaly has important physical consequences, for example, in the mass
spectrum of hadrons and in the structure of the energy-momentum tensor of
QCD.

\problems

\textbf{19.1} \emph{Fermion number nonconservation in four dimensions.}  In
this problem we study the nonconservation of fermion number in a background
field configuration with parallel electric and magnetic fields.  Consider
four-dimensional massless QED in a background field with $\mathbf{E}$ parallel
to $\mathbf{B}$.  Show that the anomaly equation leads to a violation of
fermion number, and compute the rate of fermion number violation in terms of
the field strengths.

\textbf{19.2} In the text, we derived the axial vector anomaly in
four-dimensional QED by two methods: the operator product method and the
functional integral method.  In this problem, compute the anomaly directly
from the triangle diagram, and verify that the result is
\eqref{eq:19.45}.  You will need to regulate the triangle diagram carefully,
using dimensional regularization, and keep track of the $\gamma^5$
convention.

\textbf{19.3} \emph{Anomaly coefficients for higher representations.}  In the
text we stated the anomaly coefficients $A(a) = n-4$ and $A(s) = n+4$ for
the antisymmetric and symmetric two-index tensor representations of
$SU(n)$.  Derive these results, using methods similar to those used to
compute the Casimir operators in Section~15.4.

\textbf{19.4} \emph{The trace anomaly in QCD.}  Derive the trace anomaly for
massless QCD, following the analysis of Section~19.5.  Show that the trace of
the energy-momentum tensor of QCD is proportional to
$\beta(g)\,G^a_{\mu\nu}G^{a\mu\nu}$, and discuss the physical consequences.

\textbf{19.5} In the text we showed that the anomaly factor for a chiral
gauge theory coupled to gravity is $\mathrm{tr}[t^a]$.  Show that a $U(1)$
gauge theory cannot be coupled consistently to gravity unless the $U(1)$
charge is traceless, and discuss the constraints this imposes on the matter
content of the theory.
